\documentclass[11pt,a4paper]{letter}
\usepackage[utf8]{inputenc}
\usepackage[T1]{fontenc}
\usepackage{lmodern}
\usepackage[margin=1in]{geometry}
\usepackage{hyperref}
\usepackage{xcolor}
\usepackage{amsmath}

\definecolor{primaryblue}{RGB}{0, 51, 102}

\hypersetup{
    colorlinks=true,
    linkcolor=primaryblue,
    urlcolor=primaryblue
}

\address{
    \textbf{Prof. Dr. Pedro Alves da Silva Autreto}\\
    Center of Natural and Human Sciences (CCNH)\\
    Federal University of ABC (UFABC)\\
    Santo André, SP, Brazil\\
    Email: \href{mailto:pedro.autreto@ufabc.edu.br}{pedro.autreto@ufabc.edu.br}\\
    Phone: +55 (19) 98279-4988
}

\date{August 11, 2026}

\begin{document}

\begin{letter}{
    Editor-in-Chief / Associate Editor\\
    \textit{ACS Applied Energy Materials}\\
    American Chemical Society
}

\opening{Dear Editor,}

We are pleased to submit our manuscript entitled \textbf{``Coordination-Dependent Hydrogen Adsorption on Co-, Fe-, and Ni-Functionalized Monolayer $\text{CrCl}_3$''} for consideration as a research article in \textit{ACS Applied Energy Materials}.

Developing efficient, low-cost, and stable electrocatalysts for the hydrogen evolution reaction (HER) is paramount for advancing green hydrogen technology. Two-dimensional (2D) intrinsic magnetic semiconductors, such as monolayer $\text{CrCl}_3$, offer a rich platform to couple spin degree of freedom with catalytic reactivity. However, the basal plane of pristine $\text{CrCl}_3$ is chemically inert and exhibits highly unfavorable hydrogen adsorption thermodynamics ($\Delta G_{\text{ads}} = 1.82\text{ eV}$).

In this work, using spin-polarized Density Functional Theory (DFT) calculations and \textit{ab initio} Molecular Dynamics (AIMD) simulations, we demonstrate how transition-metal ($\text{TM} = \text{Co}, \text{Fe}, \text{Ni}$) functionalization turns monolayer $\text{CrCl}_3$ into a highly promising HER catalyst by introducing localized, spin-polarized $3d$ states near the Fermi level. 

The key physical and chemical insights established in this study include:
\begin{itemize}
    \item \textbf{Coordination-Controlled Trade-off:} We discover that substrate anchoring strength does not correlate with H-binding propensity. Although Ni exhibits the strongest binding to the $\text{CrCl}_3$ host and tends to incorporate deeply into the layer at room temperature, surface-adsorbed Co achieves an optimal, near-thermoneutral adsorption free energy ($\Delta G_{\text{ads}} = 0.20\text{ eV}$).
    \item \textbf{Surface Exposure vs. Embedding:} Increasing the TM coordination from a 3-fold surface site to a 6-fold embedded site systematically degrades H adsorption performance ($\Delta G_{\text{ads}}$ up to $3.62\text{ eV}$ for embedded Ni).
    \item \textbf{Bonding Decomposition \& COHP:} Through Crystal Orbital Hamilton Population (COHP) analysis and frozen-geometry distortion-interaction partitioning, we demonstrate that higher coordination reduces the orbital accessibility of the metal $3d$ centers for H bonding while incurring large structural deformation penalties.
    \item \textbf{Limits of Standard Descriptors:} We show that conventional single-variable descriptors (such as the spin-summed $d$-band center or local magnetic moment) fail to predict adsorption trends on spin-polarized functionalized magnetic monolayers, highlighting the intrinsically multidimensional nature of active site design.
\end{itemize}

Given the strong emphasis of \textit{ACS Applied Energy Materials} on fundamental insights into material design for energy conversion applications, we believe this work will be of immediate interest to the journal's broad readership in physical chemistry, nanomaterials, and electrocatalysis.

This manuscript has not been published previously and is not currently under consideration for publication elsewhere. All authors have reviewed and approved the final manuscript for submission.

Thank you very much for considering our work.

\closing{Best wishes,}

\vspace{0.5cm}
\noindent
\textbf{Prof. Dr. Pedro Alves da Silva Autreto}\\
\textit{Corresponding Author}\\
Coordinator, Group of Electronic Structure and Atomistic Dynamics Interdisciplinary (GEEDAI)\\
Center of Natural and Human Sciences (CCNH), Federal University of ABC (UFABC)

\end{letter}
\end{document}